\section{分析骨架与核心架构}

从系统视角看，大模型推理引擎是一套同时处理执行流、状态流与控制流的运行时。模型执行与算子层决定抽象模型如何映射到具体设备能力，缓存与调度层决定请求到达过程中资源如何被持续重写，接口与运维层则决定这些内部状态能否以可升级、可观测的方式进入真实生产。对国产平台而言，很多"局部性能"问题首先暴露在模块边界处：图模式回退是否需要被调度层感知，多模态前处理应停留在接口层还是进入状态系统，平台插件究竟只负责设备发现还是还要承担环境修复与后端选择。表\ref{tab:key-tech-points} 汇总了这一分析骨架下最核心的技术问题与评价维度。

\begin{table}[t]
\centering
\small
\caption{国产推理引擎关键技术点总览}
\label{tab:key-tech-points}
\begin{tabularx}{\linewidth}{L{2.5cm}YY}
\toprule
技术主路径 & 关键技术点 & 典型评价维度 \\
\midrule
执行与通信 & 统一执行 IR、动态图与图模式切换、Prefill/Decode 阶段解耦、推测解码协同、拓扑感知通信运行时、计算--通信重叠 & MFU、吞吐、TTFT、TPOT、跨节点扩展效率、回退稳定性 \\
状态治理 & Prefix/共享状态索引、PagedAttention 与块化缓存、多级存储驻留、状态生命周期管理、跨节点状态迁移、命中感知调度 & KV 命中率、迁移流量、恢复时延、显存/主存/NVMe 占用、长上下文稳定性 \\
压缩与成本协同 & 混合精度量化、KV 动态量化、稀疏--量化协同、低比特算子路径、质量约束下的单位 token 成本优化 & 精度保持、吞吐/时延变化、状态容量变化、带宽占用、单位 token 成本 \\
统一底座与交付 & 接口协议、指标口径、trace/log/metrics、回归验证、第三方测试证据链、环境治理与部署控制面 & 可复现性、可比较性、升级回归成本、交付成熟度 \\
\bottomrule
\end{tabularx}
\end{table}

\begin{figure}[!htb]
\centering
\resizebox{0.94\linewidth}{!}{\begin{tikzpicture}[
    node distance=5mm and 6mm,
    block/.style={draw, rounded corners=2pt, align=center, minimum height=9mm, inner sep=4pt, font=\small},
    group/.style={draw, rounded corners=3pt, inner sep=5pt},
    flow/.style={-{Latex[length=2.2mm]}, semithick}
]
\node[block, fill=gray!10, minimum width=2.9cm] (ingress) {业务请求\\文本/长上下文/多模态/工具调用};
\node[block, fill=blue!8, minimum width=2.7cm, right=of ingress] (service) {服务入口与接口层\\OpenAI 兼容/鉴权/监控};
\node[block, fill=blue!8, minimum width=2.7cm, right=of service] (scheduler) {调度与运行时层\\batch 组织/SLO/bucket/回退};
\node[block, fill=green!8, minimum width=2.4cm, below=of scheduler] (cache) {KV Cache 与状态层\\前缀复用/迁移/驱逐};
\node[block, fill=green!8, minimum width=2.4cm, left=of cache] (execute) {执行层\\prefill/decode/并行策略};
\node[block, fill=green!8, minimum width=2.4cm, right=of cache] (backend) {后端接口层\\图模式/算子/能力门控};
\node[group, fit=(execute)(cache)(backend), label={[font=\small]above:运行时核心}] (runtime) {};
\node[block, fill=orange!10, minimum width=2.6cm, right=18mm of backend] (compiler) {编译器与运行时\\CANN/MindSpore/插件栈};
\node[block, fill=orange!10, minimum width=2.3cm, below=of compiler] (device) {国产硬件平台\\NPU/GPU/驱动/通信};

\draw[flow] (ingress) -- (service);
\draw[flow] (service) -- (scheduler);
\draw[flow] (scheduler) -- (backend);
\draw[flow] (scheduler) -- (execute);
\draw[flow] (scheduler) -- (cache);
\draw[flow] (execute) -- (cache);
\draw[flow] (cache) -- (backend);
\draw[flow] (backend) -- (compiler);
\draw[flow] (compiler) -- (device);
\draw[flow] (device.west) |- (backend.east);
\draw[flow, dashed] (cache.west) |- (service.south);

\node[font=\scriptsize, align=left, below=2mm of ingress] {真实压力来自多阶段请求异质性};
\node[font=\scriptsize, align=left, below=2mm of compiler] {平台差异集中暴露在编译、算子与回退边界};
\end{tikzpicture}}
\caption{国产推理引擎系统设计示意图}
\label{fig:engine-architecture-overview}
\end{figure}

\subsection{从最小推理闭环到统一分析底座}

输入文本经 tokenizer 切分、embedding 转换后进入 Transformer block 序列。对推理服务而言，关键在于 prefill 与 decode 的行为差异：Prefill 对整段提示词并行建模并建立 KV Cache，随后 decode 每步只处理增量计算和缓存更新。首 token 时延主要受前处理与 prefill 主路径影响，稳定生成速度则更多受 decode 主路径与带宽约束影响。连续批处理、chunked prefill 和阶段解耦本质上都在利用前者更接近 compute-bound、后者更接近 memory-bound 这一差异\citep{orca2023,sarathi2023,zhong2024distserve}。

\begin{figure}[!htb]
\centering
\resizebox{0.97\linewidth}{!}{\begin{tikzpicture}[
    node distance=8mm and 10mm,
    block/.style={draw=black!75, line width=0.8pt, rounded corners=2.2pt, minimum width=25mm, minimum height=10mm, align=center, font=\small, inner sep=4pt},
    note/.style={font=\scriptsize, align=center, fill=white, inner sep=1pt, text=black!80},
    flow/.style={-{Latex[length=2mm]}, semithick, draw=black!75}
]
\node[block, fill=gray!8] (input) {输入文本\\tokenizer};
\node[block, fill=blue!9, right=of input] (embed) {embedding\\位置编码};
\node[block, fill=green!10, right=of embed] (prefill) {Prefill\\建立 KV Cache};
\node[block, fill=green!10, below=12mm of prefill] (decode) {Decode\\增量生成};
\node[block, fill=gray!8, right=of decode] (output) {输出 token\\流式返回};

\draw[flow] (input) -- (embed);
\draw[flow] (embed) -- (prefill);
\draw[flow] (prefill) -- (decode);
\draw[flow] (decode) -- (output);
\draw[flow, dashed] (prefill.south west) -- ++(-9mm,0) coordinate (ttft)
    -- (ttft |- decode.west) coordinate (ttft2) -- (decode.west);
\node[note, anchor=east, xshift=-2pt, text width=2.8cm] at ($(ttft)!0.5!(ttft2)$) {首 token 时延};

\node[note, below=8mm of decode, text width=5.2cm] {稳定生成速度主要受 decode 主路径、缓存组织和带宽利用影响。};
\end{tikzpicture}}
\caption{从输入到输出的最小推理闭环}
\label{fig:minimal-inference-loop}
\end{figure}

在此基础上，可观测底座的首要作用是为多平台、多版本和多阶段优化条件下的比较提供共同语义。若缺少统一指标口径（如 TTFT、TPOT、MFU 的对齐定义），不同路线的结论就无法在同一坐标系中比较；若缺少 trace 和 metrics 基础设施，调度策略的效果就难以在长时序中被持续验证。

\begin{figure}[!htb]
\centering
\resizebox{0.85\linewidth}{!}{\definecolor{obsBlue}{HTML}{DCEEFF}
\definecolor{obsGreen}{HTML}{DDF3E4}
\definecolor{obsOrange}{HTML}{FCE8D4}
\definecolor{obsGray}{HTML}{F3F4F6}

\begin{tikzpicture}[
    node distance=11mm and 15mm,
    core/.style={circle, draw=black!70, line width=0.9pt, minimum size=22mm, align=center, font=\small\bfseries, fill=obsGray},
    petal/.style={circle, draw=black!70, line width=0.8pt, minimum size=17mm, align=center, font=\small\bfseries},
    note/.style={font=\scriptsize\sffamily, text=black!72, align=center},
    link/.style={-{Latex[length=2.2mm,width=1.8mm]}, draw=black!70, line width=0.8pt},
    guide/.style={draw=black!40, dashed, line width=0.65pt}
]

\node[core] (core) {统一\\口径};

\node[petal, fill=obsBlue, above left=12mm and 20mm of core] (api) {接口\\语义};
\node[petal, fill=obsBlue, above right=12mm and 20mm of core] (metric) {指标\\口径};
\node[petal, fill=obsGreen, below left=12mm and 20mm of core] (state) {状态\\描述};
\node[petal, fill=obsOrange, below right=12mm and 20mm of core] (trace) {观测\\证据};

\node[note, above=3mm of api] {协议 / 版本};
\node[note, above=3mm of metric] {TTFT / TPOT};
\node[note, below=3mm of state] {命中 / 迁移};
\node[note, below=3mm of trace] {trace / log};

\draw[guide] (api) -- (core);
\draw[guide] (metric) -- (core);
\draw[guide] (state) -- (core);
\draw[guide] (trace) -- (core);

\draw[link] (api.east) to[bend left=16] (metric.west);
\draw[link] (metric.south) to[bend left=16] (trace.north);
\draw[link] (trace.west) to[bend left=16] (state.east);
\draw[link] (state.north) to[bend left=16] (api.south);

\node[draw=black!60, rounded corners=2.2pt, fill=white, inner sep=4pt, align=center, font=\scriptsize\sffamily, below=21mm of core, minimum width=58mm] (outcome) {路线比较 \quad $\rightarrow$ \quad 回归验证 \quad $\rightarrow$ \quad 交付诊断};
\draw[link] (core) -- (outcome);

\node[note, right=12mm of metric, align=left] {统一后才能\\跨平台比较};
\node[note, left=12mm of state, align=right] {把指标、状态与\\日志放进同一证据链};

\end{tikzpicture}}
\caption{可观测性与能力契约层级示意图}
\label{fig:observability-contract-stack}
\end{figure}

\subsection{主路径一：执行与通信}

执行路径决定推理引擎的吞吐上限与时延下限。核心问题包括：何时使用图模式、何时回退到 eager 执行；如何在多设备间组织 tensor parallel、pipeline parallel 及 MoE expert parallel 的通信拓扑；以及如何在推测解码、chunked prefill 和阶段解耦中平衡计算与通信重叠。

对国产平台，执行路径还包含接口稳定性与维护成本。如果后端为获得局部加速在共享路径中嵌入大量平台条件分支，一旦上游框架修改模型执行接口、缓存布局或采样流程，本地后端就会承受高维护负担。FlashAttention-3 直接利用 H100 上 Tensor Core 异步执行与 FP8 支持，FlexAttention 则在保留 fused kernel 性能的同时恢复 attention 变体的可编程性\citep{shah2024flashattention3,dong2025flexattention}。这类工作说明，高性能后端的难点往往不在"有没有 kernel"，而在"如何把硬件代际能力组织成可持续复用的编译与运行时接口"。

因此，面向国产算力的设计应优先使用平台注册表、设备抽象和能力门控，把"平台差异"封装成可替换模块。算子和图编译层更需要沉淀的是一套后端能力契约：哪些算子变体可被图模式安全覆盖，哪些路径必须保留 eager 回退，哪些低比特逻辑需要运行时显式感知。对国产平台而言，很多问题并不是"后端不可用"，而是"某种能力只在某些约束下可用"。若系统无法表达这种带条件的能力边界，调度器就只能以越来越多的探测和分支维持稳定。

\begin{figure}[!htb]
\centering
\resizebox{0.88\linewidth}{!}{\begin{tikzpicture}[
    node distance=10mm and 12mm,
    block/.style={draw=black!75, line width=0.8pt, rounded corners=2.2pt, minimum width=30mm, minimum height=10.5mm, align=center, font=\small, inner sep=4pt},
    flow/.style={-{Latex[length=2mm]}, semithick, draw=black!75},
    note/.style={font=\scriptsize, align=center, fill=white, inner sep=1pt, text=black!80}
]
\node[block, fill=blue!9] (scheduler) {执行编排\\batch / 阶段切换};
\node[block, fill=green!10, right=of scheduler] (runner) {model runner\\worker / 并行策略};
\node[block, fill=orange!12, below=12mm of runner] (backend) {后端能力\\图模式 / kernel / 回退};
\node[block, fill=gray!8, right=of backend] (device) {设备与通信\\NPU / GPU / 驱动};

\draw[flow] (scheduler) -- (runner);
\draw[flow] (runner) -- (backend);
\draw[flow] (backend) -- (device);
\draw[flow, dashed] (device.north) |- (runner.east);

\node[note, below=8mm of backend, text width=6.2cm] {平台差异更适合被封装在后端能力边界内，而不是持续回流到共享调度主链。};
\end{tikzpicture}}
\caption{执行后端能力边界与回退路径}
\label{fig:execution-backend-boundary}
\end{figure}

\subsection{主路径二：状态治理与 KV Cache 组织}

KV Cache 及其元数据系统是状态治理的核心载体，直接影响并发能力、会话恢复成本与长上下文性能。PagedAttention 通过块化管理显存，降低了连续大块分配的压力，也为前缀复用、分层缓存与换入换出提供了实现基础\citep{kwon2023efficient}。对国产硬件，状态系统面临两个额外约束：不同后端对内存分配粒度与 host-device 同步开销的敏感度不同；当系统需要支持 Prefix Cache、Chunked Prefill 和跨设备迁移时，缓存元数据管理复杂度显著提升。

围绕注意力高效实现与缓存压缩已形成清晰技术主线。FlashAttention 系列从算子层逐步降低注意力计算的内存访问代价\citep{dao2022flashattention,dao2024flashattention2,shah2024flashattention3}；MInference、DuoAttention、vAttention 与 LServe 分别从动态稀疏注意力、检索头/流式头分离、虚拟内存管理和统一稀疏注意力执行改写了长上下文推理的系统组织\citep{jin2024minference,xiao2024duoattention,kwon2024vattention,yang2025lserve}；H2O、KIVI、SnapKV 与 CacheGen 则从重点击中、低比特量化、语义压缩和缓存流式传输角度缓解了 KV Cache 的显存与带宽压力\citep{zhang2023h2o,liu2024kivi,li2024snapkv,zhang2024cachegen}。

从系统演进角度看，KV Cache 正从执行层附属模块演变为独立基础设施。随着长上下文、多轮会话和多 Agent 工作流共存，缓存不仅要回答"数据放在哪"，还要回答"状态如何被计价、回收、复用和转移"。国产推理引擎在 KV Cache 设计上更适合采用"统一抽象 + 平台特化实现"的方式：统一暴露块管理、引用计数、前缀复用和驱逐策略接口，而将具体数据排布与搬运实现留给后端适配层。

\begin{figure}[!htb]
\centering
\resizebox{0.58\linewidth}{!}{\begin{tikzpicture}[
    node distance=12mm and 16mm,
    block/.style={draw=black!75, line width=0.8pt, rounded corners=2.2pt, minimum width=28mm, minimum height=10mm, align=center, font=\small, inner sep=4pt},
    flow/.style={-{Latex[length=2mm]}, semithick, draw=black!75}
]
\node[block, fill=green!10] (cache) {KV Cache\\块化 / 前缀复用};
\node[block, fill=orange!12, right=of cache] (compress) {压缩路径\\量化 / 稀疏 / 低比特};
\node[block, fill=blue!9, below=of compress] (cost) {成本结果\\显存 / 带宽 / 单位 token};
\node[block, fill=gray!8, left=of cost] (quality) {服务质量\\吞吐 / 时延 / 精度};

\draw[flow] (cache) -- (compress);
\draw[flow] (compress) -- (cost);
\draw[flow] (cost) -- (quality);
\draw[flow] (quality) -- (cache);
\end{tikzpicture}}
\caption{KV Cache、压缩与成本反馈闭环示意图}
\label{fig:kv-cost-feedback-loop}
\end{figure}

\subsection{主路径三：压缩、成本与服务质量协同}

在国产平台语境下，压缩加速不应被看作部署后期附加的"省显存技巧"，而应被视为与执行路径和状态治理并列的核心主路径。混合精度量化、KV 动态量化和稀疏--量化协同确实能直接降低显存占用与单位 token 成本，但其实际收益高度依赖后端是否存在稳定低比特路径、状态层是否允许容量变化、调度层是否能吸收 batch 形态变化。也就是说，压缩配置的效果从来不由模型精度单独决定，而由"压缩参数、执行效率、状态占用、服务质量"四者共同决定。

如图\ref{fig:kv-cost-feedback-loop} 所示，压缩路径的工程价值需要回到 KV 组织、服务质量和单位 token 成本的联动关系中才能被正确评价。QServe 通过把权重量化、KV 压缩与 serving 系统协同组织，展示了低比特部署必须在 kernel、调度和缓存之间联合设计\citep{wang2024qserve}；KIVI、SnapKV 则进一步表明，长上下文场景下的 KV 压缩只有和缓存布局与热度判断结合，才可能不显著破坏服务质量\citep{liu2024kivi,li2024snapkv}。国产推理引擎若要把"低成本"转化为系统能力，就必须让压缩配置、运行时观测与成本核算使用统一接口。

\subsection{服务层收束：调度、接口与交付闭环}

\begingroup
\sloppy
服务层需要在吞吐与时延之间平衡。典型机制包括连续批处理、请求优先级控制、SLA 感知调度和流式返回。对国产硬件，还需考虑编译图缓存、静态形状约束与后端初始化成本——调度不仅要优化"请求何时执行"，还要优化"请求如何映射到平台最擅长的执行形态"。XGrammar 说明结构化生成已需要与推理引擎共同设计\citep{dong2024xgrammar}；AI Metropolis 进一步表明，多 Agent 负载下依赖感知的乱序执行是提高并行度的关键机制\citep{xie2025aimetropolis}。

在调度策略方面，DistServe 和 Splitwise 将 Prefill 与 Decode 的资源需求进一步拆解\citep{zhong2024distserve,splitwise2024}；Llumnix 通过跨实例动态迁移请求缓解负载失衡，SOLA 则把调度目标显式对齐到 TTFT 和 TPOT 等服务级目标\citep{sun2024llumnix,hong2025sola}。TensorRT-LLM 已把多模态支持做成相对完整的运行时能力集合\citep{tensorrtllmmultimodal2026}；OmniServe 尝试把低比特量化与长上下文统一稀疏注意力纳入同一 serving 框架\citep{omniserve2025}。这些工作虽不直接针对国产算力，但为理解调度协同、量化部署和工程化取舍提供了对照。

\endgroup

状态治理推进到生产环境后，最终落到交付闭环。推理引擎的对外层不仅包括 OpenAI 兼容接口，也包括监控、日志、链路追踪和故障恢复。这一层的重要性在国产推理系统中经常被低估，但在政企部署中，系统往往首先因兼容性、可观测性或升级复杂度而被淘汰。无论是 Transformers 的模型接口、Ray Serve 的服务编排，还是 Triton Inference Server 的部署体系，都在事实上塑造了推理引擎需要兼容的工程边界\citep{transformers,rayserve,tritonserver}。

\begin{figure}[!htb]
\centering
\resizebox{0.92\linewidth}{!}{\begin{tikzpicture}[
    node distance=16mm and 18mm,
    >=Stealth,
    mainblock/.style={
        draw=black!85, 
        line width=1pt, 
        rounded corners=3pt, 
        minimum width=38mm, 
        minimum height=14mm, 
        align=center, 
        fill=gray!8, 
        inner sep=4pt
    },
    opsblock/.style={
        mainblock,
        fill=white,
        draw=black!70,
        line width=0.75pt,
        dashed
    },
    flow/.style={->, line width=1pt, draw=black!90},
    feedback/.style={->, line width=0.75pt, dashed, draw=black!70},
    boundary/.style={
        draw=black!40, 
        line width=0.5pt, 
        dashed, 
        rounded corners=4pt, 
        fill=gray!2, 
        inner sep=18pt
    },
    edgelabel/.style={
        font=\scriptsize\sffamily, 
        text=black!75, 
        fill=gray!2, 
        inner sep=2pt
    }
]

\node[mainblock] (sched) {\textbf{调度与流量整形}\\[1.5mm]\footnotesize 令牌桶 / SLO};
\node[mainblock, left=of sched] (request) {\textbf{请求进入}\\[1.5mm]\footnotesize 服务接口};
\node[mainblock, right=of sched] (exec) {\textbf{执行与返回}\\[1.5mm]\footnotesize 流式输出};

\node[opsblock, below=18mm of sched] (ops) {\textbf{观测与运维}\\[1.5mm]\footnotesize 日志采集 $\vert$ 性能监控 $\vert$ 异常回退};

\path (exec.east) ++(28mm, 0) coordinate (right_padding);

\begin{scope}[on background layer]
    \node[boundary, fit=(request) (exec) (ops) (sched) (right_padding)] (system_boundary) {};
\end{scope}

\node[draw=black!40, line width=0.5pt, fill=white, rounded corners=2pt, inner sep=4pt, 
      anchor=south east, xshift=-4mm, yshift=4mm] at (system_boundary.south east) {
    \begin{tabular}{@{}cl@{}}
    \tikz[baseline=-0.5ex]\draw[flow] (0,0) -- (0.4,0); & \scriptsize 主请求流 \\
    \tikz[baseline=-0.5ex]\draw[feedback] (0,0) -- (0.4,0); & \scriptsize 运维反馈流
    \end{tabular}
};

\draw[flow] (request) -- (sched);
\draw[flow] (sched) -- (exec);
\draw[feedback] (exec.south) |- node[pos=0.75, edgelabel, anchor=south] {监控数据采集} (ops.east);
\draw[feedback] (ops.west) -| node[pos=0.25, edgelabel, anchor=south] {异常回退 / 配置调整} (request.south);

\end{tikzpicture}}
\caption{调度、服务与运维闭环示意图}
\label{fig:service-ops-closure}
\end{figure}

综合来看，国产推理引擎的架构难点集中在模块交界处：执行层与图编译层决定可运行边界，缓存与调度层决定长上下文和并发表现，对外接口与运维层决定系统能否形成生产闭环。后文在讨论典型技术路线时，重点关注的正是不同系统如何在这些边界上做取舍。
